\documentclass[preview]{standalone}

\usepackage{amsmath}
\usepackage{amssymb}
\usepackage{stellar}
\usepackage{definitions}
\usepackage{bettelini}

\begin{document}

\id{series-of-functions-exercise}
\genpage

\section{Exercises}

\begin{snippetexercise}{cube-root-sine-function-series-exercise}{A series involving the cube root of sine}
    Consider
    \[
        f(x)=\sum_{n=1}^{\infty}
        \frac{\sqrt[3]{\sin x}}{n^2+x^2},
        \qquad x\in\realnumbers.
    \]
    Prove that \(f\) is continuous, determine where it is differentiable, and
    study its improper integrability on \(\realnumbers\).
\end{snippetexercise}

\begin{snippetsolution}{cube-root-sine-function-series-exercise-solution}{A series involving the cube root of sine}
    Write
    \[
        g(x)=\sqrt[3]{\sin x},
        \qquad
        h(x)=\sum_{n=1}^{\infty}\frac1{n^2+x^2},
    \]
    so that \(f=gh\). Since
    \[
        \left|\frac{g(x)}{n^2+x^2}\right|\leq\frac1{n^2},
        \qquad \forall x\in\realnumbers,
    \]
    the defining series for \(f\) converges totally by the Weierstrass test.
    Its terms are continuous, hence \(f\) is continuous on \(\realnumbers\).

    To study differentiability, first consider \(h\). On every compact
    interval \([-R,R]\),
    \[
        \left|\frac{-2x}{(n^2+x^2)^2}\right|
        \leq\frac{2R}{n^4}.
    \]
    Thus the series of derivatives converges locally uniformly, and
    \[
        h'(x)=-\sum_{n=1}^{\infty}\frac{2x}{(n^2+x^2)^2},
        \qquad \forall x\in\realnumbers.
    \]
    The function \(g\) is differentiable wherever \(\sin x\neq0\), so
    \(f=gh\) is differentiable on
    \(\realnumbers\difference\{k\pi\suchthat k\in\integers\}\).

    At \(x_0=k\pi\), the function \(h\) is differentiable and strictly
    positive, whereas \(g\) is not differentiable. If \(f=gh\) were
    differentiable at \(x_0\), then \(g=f/h\) would be differentiable there,
    since \(h(x_0)>0\), which is impossible. Hence \(f\) is not
    differentiable at any integer multiple of \(\pi\).

    Finally, \(g\) is periodic with mean zero, while \(h\) is positive,
    decreasing on \((0,+\infty)\), and tends to zero. Dirichlet's test for
    improper integrals therefore shows that
    \[
        \integral[0][+\infty][g(x)h(x)][x]
    \]
    converges. Since \(g\) is odd and \(h\) is even, the integral on
    \(( -\infty,0]\) also converges and the improper integral over
    \(\realnumbers\) equals zero.

    The convergence is not absolute. Indeed, for \(x\geq1\),
    \[
        h(x)\geq\sum_{1\leq n\leq x}\frac1{n^2+x^2}
        \geq\frac{\lfloor x\rfloor}{2x^2},
    \]
    so \(h(x)\) is bounded below by a positive multiple of \(1/x\) for
    large \(x\). The periodic function \(|g|\) has positive integral over
    every period, and consequently
    \(\int_1^{+\infty}|g(x)|h(x)\,dx\) diverges by comparison with the
    harmonic series.
\end{snippetsolution}

\end{document}
